\begin{proof}
La equivalencia de \emph{c)}, \emph{d)} y \emph{e)} se ha demostrado en (7.2.4).
Para mostrar que \emph{a)} y \emph{b)} son equivalentes, recordemos (7.1.9)
que afirmar que $A$ es formalmente catenario significa que, para
todo anillo cociente íntegro $B$ de $A$, $\widehat B$ es
equidimensional. Por otra parte, toda fibra del morfismo
$\operatorname{Spec}(\widehat A)\to\operatorname{Spec}(A)$ en un
punto $x\in\operatorname{Spec}(A)$ es la fibra del morfismo
$\operatorname{Spec}(\widehat B)\to\operatorname{Spec}(B)$ en el
punto genérico $x$ de $\operatorname{Spec}(B)$, donde
$B=A/\mathfrak j_x$. Que dicha fibra no tenga ningún ciclo primo
asociado inmerso equivale, por (3.3.3), a que $\widehat B$ no
tenga ideales primos asociados inmersos; esto demuestra la
equivalencia de \emph{a)} y \emph{b)}. El mismo razonamiento muestra que si
se cumple \emph{a)}, entonces, para todo anillo cociente $B$ de $A$
sin ideales primos asociados inmersos, su completado $\widehat B$
tampoco los tiene. Por otra parte, la hipótesis de que $A$ es
formalmente catenario implica que, si $B$ es equidimensional,
también lo es $\widehat B$ (7.1.8). Esto establece la última
afirmación del teorema.

Demostremos ahora que \emph{a)} implica \emph{c)}. La condición \emph{a)} implica que
$A$ es universalmente catenario (7.1.11). Demostremos además
que, para todo anillo cociente íntegro $B$ de $A$, $B^{(1)}$
es una $B$-álgebra finita. En virtud de (7.2.3) aplicado a $B$,
hay que probar lo siguiente: si $X=\operatorname{Spec}(B)$,
si $T$ es un subconjunto cerrado de $X$ de codimensión
$\geqslant2$, si $X'=\operatorname{Spec}(\widehat B)$ y si
$g:X'\to X$ es el morfismo canónico, entonces, para todo
$x'\in\operatorname{Ass}(\mathscr O_{X'})$, se tiene
$\operatorname{codim}(g^{-1}(T)\cap\overline{\{x'\}},
\overline{\{x'\}})\geqslant2$.
Pero, por hipótesis, $X'$ no tiene ningún ciclo primo asociado
inmerso. Por tanto,
$\inf(\operatorname{codim}(g^{-1}(T)\cap\overline{\{x'\}},
\overline{\{x'\}}))$ cuando $x'$ recorre
$\operatorname{Ass}(\mathscr O_{X'})$, no es sino
$\operatorname{codim}(g^{-1}(T),X')$.
Ahora bien, como $g$ es fielmente plano, se tiene (6.1.4)
\[
\operatorname{codim}(g^{-1}(T),X')=\operatorname{codim}(T,X)\geqslant2.
\]

\oldpage[IV-2]{191}
